% !LW recipe=pdflatex
\documentclass[11pt]{article}

\usepackage[margin=1truein]{geometry}
\usepackage{authblk}

\newcommand{\myQED}{}

\usepackage{graphicx}
\usepackage{mathptmx}

\usepackage{ascmac}
\usepackage{appendix}
\usepackage[linesnumbered, ruled]{algorithm2e}
\usepackage{amsmath}
\usepackage{amssymb}
\usepackage{amsthm}
\usepackage{bm,bbm}
\usepackage{booktabs}
\usepackage{cancel}
\usepackage{enumitem}
\usepackage{fancybox}
\usepackage{floatrow}
\usepackage[breaklinks]{hyperref}
\usepackage{ifthen}
\usepackage{lipsum}
\usepackage{makecell}
\usepackage{mathrsfs}
\usepackage[showonlyrefs]{mathtools}
\usepackage{multirow}
\usepackage{optidef}
\usepackage{orcidlink}
\usepackage{physics}
\usepackage{pifont}
\usepackage{subcaption}
\usepackage{subfiles}
\usepackage{thm-restate}
\usepackage{tikz}
\usepackage{xparse}

% https://tex.stackexchange.com/questions/499497/capacity-exceeded-error-while-using-cleveref-package-with-svjour3-springer-cla
\makeatletter
\let\cl@chapter\undefined
\makeatletter

\usepackage{cleveref}
\usepackage[numbers,sort&compress]{natbib}

\hypersetup{colorlinks=true,linkcolor=blue,citecolor=blue,urlcolor=blue}

\definecolor{cA}{HTML}{0072BD}
\definecolor{cB}{HTML}{EDB120}
\definecolor{cC}{HTML}{77AC30}
\definecolor{cD}{HTML}{D95319}

\newcommand{\red}[1]{\textcolor{red}{#1}}
\newcommand{\blue}[1]{\textcolor{blue}{#1}}
\newcommand{\cyan}[1]{\textcolor{cyan}{#1}}
\newcommand{\gray}[1]{\textcolor{gray}{#1}}
\newcommand{\green}[1]{\textcolor{green}{#1}}
\newcommand{\brown}[1]{\textcolor{brown}{#1}}
\newcommand{\black}[1]{\textcolor{black}{#1}}
\newcommand{\orange}[1]{\textcolor{orange}{#1}}
\newcommand{\st}{\text{ s.t. }}
\newcommand{\Img}[1]{\mathrm{Im}\qty(#1)}
\newcommand{\Ker}[1]{\mathrm{Ker}\qty(#1)}
\newcommand{\Supp}[1]{\mathrm{supp}\qty(#1)}
\newcommand{\Rank}[1]{\mathrm{rank}\qty(#1)}
\DeclarePairedDelimiter{\floor}{\lfloor}{\rfloor}
\DeclarePairedDelimiter{\ceil}{\lceil}{\rceil}
% C++ (https://tex.stackexchange.com/questions/4302/prettiest-way-to-typeset-c-cplusplus)
\newcommand{\Cpp}{C\nolinebreak[4]\hspace{-.05em}\raisebox{.4ex}{\relsize{-3}{\textbf{++}}}}
% https://tex.stackexchange.com/questions/28836/typesetting-the-define-equals-symbol
\newcommand{\defeq}{\coloneqq}
\newcommand{\eqdef}{\eqqcolon}
% https://tex.stackexchange.com/questions/5502/how-to-get-a-mid-binary-relation-that-grows
\newcommand{\relmiddle}[1]{\mathrel{}\middle#1\mathrel{}}

% https://tex.stackexchange.com/questions/564216/newcommand-for-each-letter
\ExplSyntaxOn
\NewDocumentCommand{\definealphabet}{mmmm}{
    \int_step_inline:nnn{`#3}{`#4}{
        \cs_new_protected:cpx{#1 \char_generate:nn{##1}{11}}{
            \exp_not:N #2{\char_generate:nn{##1}{11}}}}}
\ExplSyntaxOff

% cleveref fails to accurately capture pseudo-code line numbers. We define a custom reference command for lines in algorithms.
\newcommand{\refLine}[1]{line~\ref{#1}}
\newcommand{\refLines}[2]{lines~\ref{#1}--\ref{#2}}

\definealphabet{bb}{\mathbb}{A}{Z}
\definealphabet{rm}{\mathrm}{A}{Z}
\definealphabet{cal}{\mathcal}{A}{Z}
\definealphabet{frak}{\mathfrak}{a}{z}
% \definealphabet{scr}{\mathscr}{A}{Z}
% \definealphabet{frak}{\mathfrak}{A}{Z}

\newtheorem{theorem}{Theorem}
\newtheorem{proposition}{Proposition}
\newtheorem{lemma}{Lemma}
\newtheorem{corollary}{Corollary}
\newtheorem{assumption}{Assumption}

\theoremstyle{definition}
\newtheorem{definition}{Definition}
\newtheorem{example}{Example}

\theoremstyle{remark}
\newtheorem{remark}{Remark}

% cleverefの設定
\crefname{theorem}{Theorem}{Theorems}
\crefname{proposition}{Proposition}{Propositions}
\crefname{lemma}{Lemma}{Lemmas}
\crefname{assumption}{Assumption}{Assumptions}
\crefname{definition}{Definition}{Definitions}
\crefname{corollary}{Corollary}{Corollaries}
\crefname{remark}{Remark}{Remarks}
\crefname{example}{Example}{Examples}
\crefname{figure}{Figure}{Figures}
\crefname{table}{Table}{Tables}
\crefname{algorithm}{Algorithm}{Algorithms}

% https://qiita.com/rityo_masu/items/efd44bc8f9229e014237
\allowdisplaybreaks[4]

\newcommand{\feps}{\epsilon_\mathrm{f}}
\newcommand{\gtol}{\epsilon_{\mathrm{gtol}}}

\begin{document}

\title{Practical Regularized Quasi-Newton Methods\\
    with Inexact Function Values}

\author[1]{Hiroki Hamaguchi\,\orcidlink{0009-0005-7348-1356}\footnote{E-mail: \href{mailto:hirokihamaguchi0331@gmail.com}{\nolinkurl{hirokihamaguchi0331@gmail.com}}}}
\author[1]{Naoki Marumo\,\orcidlink{0000-0002-7372-4275}}
\author[1,2]{Akiko Takeda\,\orcidlink{0000-0002-8846-4496}}

\affil[1]{Graduate School of Information Science and Technology, The University of Tokyo, Tokyo, Japan}
\affil[2]{Center for Advanced Intelligence Project, RIKEN, Tokyo, Japan}

\maketitle

\begin{abstract}
    Many practical optimization problems involve objective function values that are corrupted by \MPCchanged{bounded, nonvanishing} numerical errors.
\MPCchanged{We} propose a noise-tolerant regularized quasi-Newton method \MPCchanged{that uses a relaxed Armijo line search when observed function values certify progress and switches to OFFO regularization otherwise}.
\MPCchanged{With exact gradients and uniform spectral bounds on the quasi-Newton matrices, we prove an order-optimal \(\order{1/\varepsilon^2}\) complexity bound under bounded relative-or-absolute errors in function values.}
\MPCchanged{Although the method uses persistently noisy function values when they certify progress, it retains the optimal worst-case dependence on the stationarity tolerance for the noise-free $L$-smooth nonconvex setting.}
\MPCchanged{The exact-gradient qualification is essential to the stated theorem.}
\MPCchanged{Experiments with perturbed gradients are reported separately as deliberately adverse robustness tests beyond its scope.}
\MPCchanged{Experiments on CUTEst include function-value-only noise, joint function-and-gradient noise, and 64-, 32-, and 16-bit arithmetic.}
\MPCchanged{They evaluate both robustness and the sensitivity to the supplied function-error bound.}

\end{abstract}

\section{Introduction}\label{sec:introduction}

\subfile{introduction.tex}

The remainder of this paper is organized as follows.
In \cref{sec:preliminaries}, we review relevant background and introduce the problem setting.
In \cref{sec:algorithm}, we present the proposed algorithm.
In \cref{sec:analysis}, we establish a global convergence property of the proposed algorithm.
In \cref{sec:proposed_choice}, we discuss practical choices of algorithmic variables and parameters.
In \cref{sec:experiments}, we report numerical results on the CUTEst benchmark collection.
Finally, in \cref{sec:conclusion}, we conclude the paper with a summary and future research directions.

Our implementation is publicly available on GitHub%
\footnote{\label{fn:ourGitHub}\url{https://github.com/HirokiHamaguchi/qnlab}}.

% arXiv-specific widths for Figure~\ref{fig:results_noise} in main_contents.tex
\providecommand{\resultsNoiseMainWidth}{0.7}
\providecommand{\resultsNoiseLegendWidth}{0.7}
\providecommand{\resultsAllPrecisionsSubfigWidth}{0.25}

\subfile{main_contents.tex}

\section*{Acknowledgements}
This work was partially supported by JSPS KAKENHI (23H03351, 24K23853) and JST CREST (JPMJCR24Q2).

% % BibTeX users please use one of
% % \bibliographystyle{spbasic}      % basic style, author-year citations
% \bibliographystyle{spmpsci}      % mathematics and physical sciences
% %\bibliographystyle{spphys}       % APS-like style for physics
% \bibliography{../L-BFGS.bib}
% !LW recipe=pdflatex
\documentclass[11pt]{article}

\usepackage[margin=1truein]{geometry}
\usepackage{authblk}

\newcommand{\myQED}{}

\usepackage{graphicx}
\usepackage{mathptmx}

\usepackage{ascmac}
\usepackage{appendix}
\usepackage[linesnumbered, ruled]{algorithm2e}
\usepackage{amsmath}
\usepackage{amssymb}
\usepackage{amsthm}
\usepackage{bm,bbm}
\usepackage{booktabs}
\usepackage{cancel}
\usepackage{enumitem}
\usepackage{fancybox}
\usepackage{floatrow}
\usepackage[breaklinks]{hyperref}
\usepackage{ifthen}
\usepackage{lipsum}
\usepackage{makecell}
\usepackage{mathrsfs}
\usepackage[showonlyrefs]{mathtools}
\usepackage{multirow}
\usepackage{optidef}
\usepackage{orcidlink}
\usepackage{physics}
\usepackage{pifont}
\usepackage{subcaption}
\usepackage{subfiles}
\usepackage{thm-restate}
\usepackage{tikz}
\usepackage{xparse}

% https://tex.stackexchange.com/questions/499497/capacity-exceeded-error-while-using-cleveref-package-with-svjour3-springer-cla
\makeatletter
\let\cl@chapter\undefined
\makeatletter

\usepackage{cleveref}
\usepackage[numbers,sort&compress]{natbib}

\hypersetup{colorlinks=true,linkcolor=blue,citecolor=blue,urlcolor=blue}

\definecolor{cA}{HTML}{0072BD}
\definecolor{cB}{HTML}{EDB120}
\definecolor{cC}{HTML}{77AC30}
\definecolor{cD}{HTML}{D95319}

\newcommand{\red}[1]{\textcolor{red}{#1}}
\newcommand{\blue}[1]{\textcolor{blue}{#1}}
\newcommand{\cyan}[1]{\textcolor{cyan}{#1}}
\newcommand{\gray}[1]{\textcolor{gray}{#1}}
\newcommand{\green}[1]{\textcolor{green}{#1}}
\newcommand{\brown}[1]{\textcolor{brown}{#1}}
\newcommand{\black}[1]{\textcolor{black}{#1}}
\newcommand{\orange}[1]{\textcolor{orange}{#1}}
\newcommand{\st}{\text{ s.t. }}
\newcommand{\Img}[1]{\mathrm{Im}\qty(#1)}
\newcommand{\Ker}[1]{\mathrm{Ker}\qty(#1)}
\newcommand{\Supp}[1]{\mathrm{supp}\qty(#1)}
\newcommand{\Rank}[1]{\mathrm{rank}\qty(#1)}
\DeclarePairedDelimiter{\floor}{\lfloor}{\rfloor}
\DeclarePairedDelimiter{\ceil}{\lceil}{\rceil}
% C++ (https://tex.stackexchange.com/questions/4302/prettiest-way-to-typeset-c-cplusplus)
\newcommand{\Cpp}{C\nolinebreak[4]\hspace{-.05em}\raisebox{.4ex}{\relsize{-3}{\textbf{++}}}}
% https://tex.stackexchange.com/questions/28836/typesetting-the-define-equals-symbol
\newcommand{\defeq}{\coloneqq}
\newcommand{\eqdef}{\eqqcolon}
% https://tex.stackexchange.com/questions/5502/how-to-get-a-mid-binary-relation-that-grows
\newcommand{\relmiddle}[1]{\mathrel{}\middle#1\mathrel{}}

% https://tex.stackexchange.com/questions/564216/newcommand-for-each-letter
\ExplSyntaxOn
\NewDocumentCommand{\definealphabet}{mmmm}{
    \int_step_inline:nnn{`#3}{`#4}{
        \cs_new_protected:cpx{#1 \char_generate:nn{##1}{11}}{
            \exp_not:N #2{\char_generate:nn{##1}{11}}}}}
\ExplSyntaxOff

% cleveref fails to accurately capture pseudo-code line numbers. We define a custom reference command for lines in algorithms.
\newcommand{\refLine}[1]{line~\ref{#1}}
\newcommand{\refLines}[2]{lines~\ref{#1}--\ref{#2}}

\definealphabet{bb}{\mathbb}{A}{Z}
\definealphabet{rm}{\mathrm}{A}{Z}
\definealphabet{cal}{\mathcal}{A}{Z}
\definealphabet{frak}{\mathfrak}{a}{z}
% \definealphabet{scr}{\mathscr}{A}{Z}
% \definealphabet{frak}{\mathfrak}{A}{Z}

\newtheorem{theorem}{Theorem}
\newtheorem{proposition}{Proposition}
\newtheorem{lemma}{Lemma}
\newtheorem{corollary}{Corollary}
\newtheorem{assumption}{Assumption}

\theoremstyle{definition}
\newtheorem{definition}{Definition}
\newtheorem{example}{Example}

\theoremstyle{remark}
\newtheorem{remark}{Remark}

% cleverefの設定
\crefname{theorem}{Theorem}{Theorems}
\crefname{proposition}{Proposition}{Propositions}
\crefname{lemma}{Lemma}{Lemmas}
\crefname{assumption}{Assumption}{Assumptions}
\crefname{definition}{Definition}{Definitions}
\crefname{corollary}{Corollary}{Corollaries}
\crefname{remark}{Remark}{Remarks}
\crefname{example}{Example}{Examples}
\crefname{figure}{Figure}{Figures}
\crefname{table}{Table}{Tables}
\crefname{algorithm}{Algorithm}{Algorithms}

% https://qiita.com/rityo_masu/items/efd44bc8f9229e014237
\allowdisplaybreaks[4]

\newcommand{\feps}{\epsilon_\mathrm{f}}
\newcommand{\gtol}{\epsilon_{\mathrm{gtol}}}

\begin{document}

\title{Practical Regularized Quasi-Newton Methods\\
    with Inexact Function Values}

\author[1]{Hiroki Hamaguchi\,\orcidlink{0009-0005-7348-1356}\footnote{E-mail: \href{mailto:hirokihamaguchi0331@gmail.com}{\nolinkurl{hirokihamaguchi0331@gmail.com}}}}
\author[1]{Naoki Marumo\,\orcidlink{0000-0002-7372-4275}}
\author[1,2]{Akiko Takeda\,\orcidlink{0000-0002-8846-4496}}

\affil[1]{Graduate School of Information Science and Technology, The University of Tokyo, Tokyo, Japan}
\affil[2]{Center for Advanced Intelligence Project, RIKEN, Tokyo, Japan}

\maketitle

\begin{abstract}
    Many practical optimization problems involve objective function values that are corrupted by \MPCchanged{bounded, nonvanishing} numerical errors.
\MPCchanged{We} propose a noise-tolerant regularized quasi-Newton method \MPCchanged{that uses a relaxed Armijo line search when observed function values certify progress and switches to OFFO regularization otherwise}.
\MPCchanged{With exact gradients and uniform spectral bounds on the quasi-Newton matrices, we prove an order-optimal \(\order{1/\varepsilon^2}\) complexity bound under bounded relative-or-absolute errors in function values.}
\MPCchanged{Although the method uses persistently noisy function values when they certify progress, it retains the optimal worst-case dependence on the stationarity tolerance for the noise-free $L$-smooth nonconvex setting.}
\MPCchanged{The exact-gradient qualification is essential to the stated theorem.}
\MPCchanged{Experiments with perturbed gradients are reported separately as deliberately adverse robustness tests beyond its scope.}
\MPCchanged{Experiments on CUTEst include function-value-only noise, joint function-and-gradient noise, and 64-, 32-, and 16-bit arithmetic.}
\MPCchanged{They evaluate both robustness and the sensitivity to the supplied function-error bound.}

\end{abstract}

\section{Introduction}\label{sec:introduction}

\subfile{introduction.tex}

The remainder of this paper is organized as follows.
In \cref{sec:preliminaries}, we review relevant background and introduce the problem setting.
In \cref{sec:algorithm}, we present the proposed algorithm.
In \cref{sec:analysis}, we establish a global convergence property of the proposed algorithm.
In \cref{sec:proposed_choice}, we discuss practical choices of algorithmic variables and parameters.
In \cref{sec:experiments}, we report numerical results on the CUTEst benchmark collection.
Finally, in \cref{sec:conclusion}, we conclude the paper with a summary and future research directions.

Our implementation is publicly available on GitHub%
\footnote{\label{fn:ourGitHub}\url{https://github.com/HirokiHamaguchi/qnlab}}.

% arXiv-specific widths for Figure~\ref{fig:results_noise} in main_contents.tex
\providecommand{\resultsNoiseMainWidth}{0.7}
\providecommand{\resultsNoiseLegendWidth}{0.7}
\providecommand{\resultsAllPrecisionsSubfigWidth}{0.25}

\subfile{main_contents.tex}

\section*{Acknowledgements}
This work was partially supported by JSPS KAKENHI (23H03351, 24K23853) and JST CREST (JPMJCR24Q2).

% % BibTeX users please use one of
% % \bibliographystyle{spbasic}      % basic style, author-year citations
% \bibliographystyle{spmpsci}      % mathematics and physical sciences
% %\bibliographystyle{spphys}       % APS-like style for physics
% \bibliography{../L-BFGS.bib}
% !LW recipe=pdflatex
\documentclass[11pt]{article}

\usepackage[margin=1truein]{geometry}
\usepackage{authblk}

\newcommand{\myQED}{}

\usepackage{graphicx}
\usepackage{mathptmx}

\usepackage{ascmac}
\usepackage{appendix}
\usepackage[linesnumbered, ruled]{algorithm2e}
\usepackage{amsmath}
\usepackage{amssymb}
\usepackage{amsthm}
\usepackage{bm,bbm}
\usepackage{booktabs}
\usepackage{cancel}
\usepackage{enumitem}
\usepackage{fancybox}
\usepackage{floatrow}
\usepackage[breaklinks]{hyperref}
\usepackage{ifthen}
\usepackage{lipsum}
\usepackage{makecell}
\usepackage{mathrsfs}
\usepackage[showonlyrefs]{mathtools}
\usepackage{multirow}
\usepackage{optidef}
\usepackage{orcidlink}
\usepackage{physics}
\usepackage{pifont}
\usepackage{subcaption}
\usepackage{subfiles}
\usepackage{thm-restate}
\usepackage{tikz}
\usepackage{xparse}

% https://tex.stackexchange.com/questions/499497/capacity-exceeded-error-while-using-cleveref-package-with-svjour3-springer-cla
\makeatletter
\let\cl@chapter\undefined
\makeatletter

\usepackage{cleveref}
\usepackage[numbers,sort&compress]{natbib}

\hypersetup{colorlinks=true,linkcolor=blue,citecolor=blue,urlcolor=blue}

\definecolor{cA}{HTML}{0072BD}
\definecolor{cB}{HTML}{EDB120}
\definecolor{cC}{HTML}{77AC30}
\definecolor{cD}{HTML}{D95319}

\newcommand{\red}[1]{\textcolor{red}{#1}}
\newcommand{\blue}[1]{\textcolor{blue}{#1}}
\newcommand{\cyan}[1]{\textcolor{cyan}{#1}}
\newcommand{\gray}[1]{\textcolor{gray}{#1}}
\newcommand{\green}[1]{\textcolor{green}{#1}}
\newcommand{\brown}[1]{\textcolor{brown}{#1}}
\newcommand{\black}[1]{\textcolor{black}{#1}}
\newcommand{\orange}[1]{\textcolor{orange}{#1}}
\newcommand{\st}{\text{ s.t. }}
\newcommand{\Img}[1]{\mathrm{Im}\qty(#1)}
\newcommand{\Ker}[1]{\mathrm{Ker}\qty(#1)}
\newcommand{\Supp}[1]{\mathrm{supp}\qty(#1)}
\newcommand{\Rank}[1]{\mathrm{rank}\qty(#1)}
\DeclarePairedDelimiter{\floor}{\lfloor}{\rfloor}
\DeclarePairedDelimiter{\ceil}{\lceil}{\rceil}
% C++ (https://tex.stackexchange.com/questions/4302/prettiest-way-to-typeset-c-cplusplus)
\newcommand{\Cpp}{C\nolinebreak[4]\hspace{-.05em}\raisebox{.4ex}{\relsize{-3}{\textbf{++}}}}
% https://tex.stackexchange.com/questions/28836/typesetting-the-define-equals-symbol
\newcommand{\defeq}{\coloneqq}
\newcommand{\eqdef}{\eqqcolon}
% https://tex.stackexchange.com/questions/5502/how-to-get-a-mid-binary-relation-that-grows
\newcommand{\relmiddle}[1]{\mathrel{}\middle#1\mathrel{}}

% https://tex.stackexchange.com/questions/564216/newcommand-for-each-letter
\ExplSyntaxOn
\NewDocumentCommand{\definealphabet}{mmmm}{
    \int_step_inline:nnn{`#3}{`#4}{
        \cs_new_protected:cpx{#1 \char_generate:nn{##1}{11}}{
            \exp_not:N #2{\char_generate:nn{##1}{11}}}}}
\ExplSyntaxOff

% cleveref fails to accurately capture pseudo-code line numbers. We define a custom reference command for lines in algorithms.
\newcommand{\refLine}[1]{line~\ref{#1}}
\newcommand{\refLines}[2]{lines~\ref{#1}--\ref{#2}}

\definealphabet{bb}{\mathbb}{A}{Z}
\definealphabet{rm}{\mathrm}{A}{Z}
\definealphabet{cal}{\mathcal}{A}{Z}
\definealphabet{frak}{\mathfrak}{a}{z}
% \definealphabet{scr}{\mathscr}{A}{Z}
% \definealphabet{frak}{\mathfrak}{A}{Z}

\newtheorem{theorem}{Theorem}
\newtheorem{proposition}{Proposition}
\newtheorem{lemma}{Lemma}
\newtheorem{corollary}{Corollary}
\newtheorem{assumption}{Assumption}

\theoremstyle{definition}
\newtheorem{definition}{Definition}
\newtheorem{example}{Example}

\theoremstyle{remark}
\newtheorem{remark}{Remark}

% cleverefの設定
\crefname{theorem}{Theorem}{Theorems}
\crefname{proposition}{Proposition}{Propositions}
\crefname{lemma}{Lemma}{Lemmas}
\crefname{assumption}{Assumption}{Assumptions}
\crefname{definition}{Definition}{Definitions}
\crefname{corollary}{Corollary}{Corollaries}
\crefname{remark}{Remark}{Remarks}
\crefname{example}{Example}{Examples}
\crefname{figure}{Figure}{Figures}
\crefname{table}{Table}{Tables}
\crefname{algorithm}{Algorithm}{Algorithms}

% https://qiita.com/rityo_masu/items/efd44bc8f9229e014237
\allowdisplaybreaks[4]

\newcommand{\feps}{\epsilon_\mathrm{f}}
\newcommand{\gtol}{\epsilon_{\mathrm{gtol}}}

\begin{document}

\title{Practical Regularized Quasi-Newton Methods\\
    with Inexact Function Values}

\author[1]{Hiroki Hamaguchi\,\orcidlink{0009-0005-7348-1356}\footnote{E-mail: \href{mailto:hirokihamaguchi0331@gmail.com}{\nolinkurl{hirokihamaguchi0331@gmail.com}}}}
\author[1]{Naoki Marumo\,\orcidlink{0000-0002-7372-4275}}
\author[1,2]{Akiko Takeda\,\orcidlink{0000-0002-8846-4496}}

\affil[1]{Graduate School of Information Science and Technology, The University of Tokyo, Tokyo, Japan}
\affil[2]{Center for Advanced Intelligence Project, RIKEN, Tokyo, Japan}

\maketitle

\begin{abstract}
    Many practical optimization problems involve objective function values that are corrupted by \MPCchanged{bounded, nonvanishing} numerical errors.
\MPCchanged{We} propose a noise-tolerant regularized quasi-Newton method \MPCchanged{that uses a relaxed Armijo line search when observed function values certify progress and switches to OFFO regularization otherwise}.
\MPCchanged{With exact gradients and uniform spectral bounds on the quasi-Newton matrices, we prove an order-optimal \(\order{1/\varepsilon^2}\) complexity bound under bounded relative-or-absolute errors in function values.}
\MPCchanged{Although the method uses persistently noisy function values when they certify progress, it retains the optimal worst-case dependence on the stationarity tolerance for the noise-free $L$-smooth nonconvex setting.}
\MPCchanged{The exact-gradient qualification is essential to the stated theorem.}
\MPCchanged{Experiments with perturbed gradients are reported separately as deliberately adverse robustness tests beyond its scope.}
\MPCchanged{Experiments on CUTEst include function-value-only noise, joint function-and-gradient noise, and 64-, 32-, and 16-bit arithmetic.}
\MPCchanged{They evaluate both robustness and the sensitivity to the supplied function-error bound.}

\end{abstract}

\section{Introduction}\label{sec:introduction}

\subfile{introduction.tex}

The remainder of this paper is organized as follows.
In \cref{sec:preliminaries}, we review relevant background and introduce the problem setting.
In \cref{sec:algorithm}, we present the proposed algorithm.
In \cref{sec:analysis}, we establish a global convergence property of the proposed algorithm.
In \cref{sec:proposed_choice}, we discuss practical choices of algorithmic variables and parameters.
In \cref{sec:experiments}, we report numerical results on the CUTEst benchmark collection.
Finally, in \cref{sec:conclusion}, we conclude the paper with a summary and future research directions.

Our implementation is publicly available on GitHub%
\footnote{\label{fn:ourGitHub}\url{https://github.com/HirokiHamaguchi/qnlab}}.

% arXiv-specific widths for Figure~\ref{fig:results_noise} in main_contents.tex
\providecommand{\resultsNoiseMainWidth}{0.7}
\providecommand{\resultsNoiseLegendWidth}{0.7}
\providecommand{\resultsAllPrecisionsSubfigWidth}{0.25}

\subfile{main_contents.tex}

\section*{Acknowledgements}
This work was partially supported by JSPS KAKENHI (23H03351, 24K23853) and JST CREST (JPMJCR24Q2).

% % BibTeX users please use one of
% % \bibliographystyle{spbasic}      % basic style, author-year citations
% \bibliographystyle{spmpsci}      % mathematics and physical sciences
% %\bibliographystyle{spphys}       % APS-like style for physics
% \bibliography{../L-BFGS.bib}
% !LW recipe=pdflatex
\documentclass[11pt]{article}

\usepackage[margin=1truein]{geometry}
\usepackage{authblk}

\newcommand{\myQED}{}

\usepackage{graphicx}
\usepackage{mathptmx}

\usepackage{ascmac}
\usepackage{appendix}
\usepackage[linesnumbered, ruled]{algorithm2e}
\usepackage{amsmath}
\usepackage{amssymb}
\usepackage{amsthm}
\usepackage{bm,bbm}
\usepackage{booktabs}
\usepackage{cancel}
\usepackage{enumitem}
\usepackage{fancybox}
\usepackage{floatrow}
\usepackage[breaklinks]{hyperref}
\usepackage{ifthen}
\usepackage{lipsum}
\usepackage{makecell}
\usepackage{mathrsfs}
\usepackage[showonlyrefs]{mathtools}
\usepackage{multirow}
\usepackage{optidef}
\usepackage{orcidlink}
\usepackage{physics}
\usepackage{pifont}
\usepackage{subcaption}
\usepackage{subfiles}
\usepackage{thm-restate}
\usepackage{tikz}
\usepackage{xparse}

% https://tex.stackexchange.com/questions/499497/capacity-exceeded-error-while-using-cleveref-package-with-svjour3-springer-cla
\makeatletter
\let\cl@chapter\undefined
\makeatletter

\usepackage{cleveref}
\usepackage[numbers,sort&compress]{natbib}

\hypersetup{colorlinks=true,linkcolor=blue,citecolor=blue,urlcolor=blue}

\definecolor{cA}{HTML}{0072BD}
\definecolor{cB}{HTML}{EDB120}
\definecolor{cC}{HTML}{77AC30}
\definecolor{cD}{HTML}{D95319}

\newcommand{\red}[1]{\textcolor{red}{#1}}
\newcommand{\blue}[1]{\textcolor{blue}{#1}}
\newcommand{\cyan}[1]{\textcolor{cyan}{#1}}
\newcommand{\gray}[1]{\textcolor{gray}{#1}}
\newcommand{\green}[1]{\textcolor{green}{#1}}
\newcommand{\brown}[1]{\textcolor{brown}{#1}}
\newcommand{\black}[1]{\textcolor{black}{#1}}
\newcommand{\orange}[1]{\textcolor{orange}{#1}}
\newcommand{\st}{\text{ s.t. }}
\newcommand{\Img}[1]{\mathrm{Im}\qty(#1)}
\newcommand{\Ker}[1]{\mathrm{Ker}\qty(#1)}
\newcommand{\Supp}[1]{\mathrm{supp}\qty(#1)}
\newcommand{\Rank}[1]{\mathrm{rank}\qty(#1)}
\DeclarePairedDelimiter{\floor}{\lfloor}{\rfloor}
\DeclarePairedDelimiter{\ceil}{\lceil}{\rceil}
% C++ (https://tex.stackexchange.com/questions/4302/prettiest-way-to-typeset-c-cplusplus)
\newcommand{\Cpp}{C\nolinebreak[4]\hspace{-.05em}\raisebox{.4ex}{\relsize{-3}{\textbf{++}}}}
% https://tex.stackexchange.com/questions/28836/typesetting-the-define-equals-symbol
\newcommand{\defeq}{\coloneqq}
\newcommand{\eqdef}{\eqqcolon}
% https://tex.stackexchange.com/questions/5502/how-to-get-a-mid-binary-relation-that-grows
\newcommand{\relmiddle}[1]{\mathrel{}\middle#1\mathrel{}}

% https://tex.stackexchange.com/questions/564216/newcommand-for-each-letter
\ExplSyntaxOn
\NewDocumentCommand{\definealphabet}{mmmm}{
    \int_step_inline:nnn{`#3}{`#4}{
        \cs_new_protected:cpx{#1 \char_generate:nn{##1}{11}}{
            \exp_not:N #2{\char_generate:nn{##1}{11}}}}}
\ExplSyntaxOff

% cleveref fails to accurately capture pseudo-code line numbers. We define a custom reference command for lines in algorithms.
\newcommand{\refLine}[1]{line~\ref{#1}}
\newcommand{\refLines}[2]{lines~\ref{#1}--\ref{#2}}

\definealphabet{bb}{\mathbb}{A}{Z}
\definealphabet{rm}{\mathrm}{A}{Z}
\definealphabet{cal}{\mathcal}{A}{Z}
\definealphabet{frak}{\mathfrak}{a}{z}
% \definealphabet{scr}{\mathscr}{A}{Z}
% \definealphabet{frak}{\mathfrak}{A}{Z}

\newtheorem{theorem}{Theorem}
\newtheorem{proposition}{Proposition}
\newtheorem{lemma}{Lemma}
\newtheorem{corollary}{Corollary}
\newtheorem{assumption}{Assumption}

\theoremstyle{definition}
\newtheorem{definition}{Definition}
\newtheorem{example}{Example}

\theoremstyle{remark}
\newtheorem{remark}{Remark}

% cleverefの設定
\crefname{theorem}{Theorem}{Theorems}
\crefname{proposition}{Proposition}{Propositions}
\crefname{lemma}{Lemma}{Lemmas}
\crefname{assumption}{Assumption}{Assumptions}
\crefname{definition}{Definition}{Definitions}
\crefname{corollary}{Corollary}{Corollaries}
\crefname{remark}{Remark}{Remarks}
\crefname{example}{Example}{Examples}
\crefname{figure}{Figure}{Figures}
\crefname{table}{Table}{Tables}
\crefname{algorithm}{Algorithm}{Algorithms}

% https://qiita.com/rityo_masu/items/efd44bc8f9229e014237
\allowdisplaybreaks[4]

\newcommand{\feps}{\epsilon_\mathrm{f}}
\newcommand{\gtol}{\epsilon_{\mathrm{gtol}}}

\begin{document}

\title{Practical Regularized Quasi-Newton Methods\\
    with Inexact Function Values}

\author[1]{Hiroki Hamaguchi\,\orcidlink{0009-0005-7348-1356}\footnote{E-mail: \href{mailto:hirokihamaguchi0331@gmail.com}{\nolinkurl{hirokihamaguchi0331@gmail.com}}}}
\author[1]{Naoki Marumo\,\orcidlink{0000-0002-7372-4275}}
\author[1,2]{Akiko Takeda\,\orcidlink{0000-0002-8846-4496}}

\affil[1]{Graduate School of Information Science and Technology, The University of Tokyo, Tokyo, Japan}
\affil[2]{Center for Advanced Intelligence Project, RIKEN, Tokyo, Japan}

\maketitle

\begin{abstract}
    \input{abstract.tex}
\end{abstract}

\section{Introduction}\label{sec:introduction}

\subfile{introduction.tex}

The remainder of this paper is organized as follows.
In \cref{sec:preliminaries}, we review relevant background and introduce the problem setting.
In \cref{sec:algorithm}, we present the proposed algorithm.
In \cref{sec:analysis}, we establish a global convergence property of the proposed algorithm.
In \cref{sec:proposed_choice}, we discuss practical choices of algorithmic variables and parameters.
In \cref{sec:experiments}, we report numerical results on the CUTEst benchmark collection.
Finally, in \cref{sec:conclusion}, we conclude the paper with a summary and future research directions.

Our implementation is publicly available on GitHub%
\footnote{\label{fn:ourGitHub}\url{https://github.com/HirokiHamaguchi/qnlab}}.

% arXiv-specific widths for Figure~\ref{fig:results_noise} in main_contents.tex
\providecommand{\resultsNoiseMainWidth}{0.7}
\providecommand{\resultsNoiseLegendWidth}{0.7}
\providecommand{\resultsAllPrecisionsSubfigWidth}{0.25}

\subfile{main_contents.tex}

\section*{Acknowledgements}
This work was partially supported by JSPS KAKENHI (23H03351, 24K23853) and JST CREST (JPMJCR24Q2).

% % BibTeX users please use one of
% % \bibliographystyle{spbasic}      % basic style, author-year citations
% \bibliographystyle{spmpsci}      % mathematics and physical sciences
% %\bibliographystyle{spphys}       % APS-like style for physics
% \bibliography{../L-BFGS.bib}
\input{arXiv.bbl}

\begin{appendices}
    \crefalias{section}{appendix}
    \input{appendix.tex}
\end{appendices}

\end{document}


\begin{appendices}
    \crefalias{section}{appendix}
    \section{Proof of \texorpdfstring{\cref{lem:sum_norm_g_over_mu_geq_2sqrt_leq_log}}{Lemma \ref{lem:sum_norm_g_over_mu_geq_2sqrt_leq_log}}}\label{sec:proof_of_sum_norm_g_over_mu_geq_2sqrt_leq_log}

The proof of \cref{lem:sum_norm_g_over_mu_geq_2sqrt_leq_log} is \MPCchanged{similar to previous arguments~\citep[Lemma~3.2]{wardAdaGradStepsizesSharp2020} and \citep[Lemma~3.1]{Gratton08022024}}.
\begin{proof}
    Recall that \cref{eq:offo_based_update} defines $\mu_k$ as $\theta_k \sqrt{\varsigma + \sum_{j \in K^+, \, j \leq k} \norm{g_j}^2}$, and \refLine{line:update_theta} of \cref{alg:regularized_qn} asserts $\theta_k \in [\theta_{\min}, \theta_{\max}]$.
    Let $k_i \in K^+$ be the $i$-th element of $K^+$.
    \MPCchanged{Define $S_{-1} \coloneqq \varsigma$ and $S_i \coloneqq \varsigma + \sum_{\ell=0}^{i} \norm{g_{k_\ell}}^2$ for $i \geq 0$.}
    For any $0 < x < y$, $x/y \le -\ln(1-x/y) = \ln y - \ln(y-x)$ holds.
    Thus, we have
    \begin{align*}
        \frac{\norm{g_{k_i}}^2}{\mu_{k_i}^2}
         & \MPCchanged{\leq \frac{1}{\theta_{\min}^2} \frac{\norm{g_{k_i}}^2}{S_i}} &  & (\text{definition}) \\
         & \MPCchanged{\leq \frac{1}{\theta_{\min}^2} \qty(\ln S_i - \ln S_{i-1})}.  &  & (x/y \leq \ln y - \ln(y-x))
    \end{align*}
    Similarly, $x/\sqrt{y} \ge x/(\sqrt{y}+\sqrt{y-x}) = \sqrt{y}-\sqrt{y-x}$ holds.
    Thus, we have
    \begin{align*}
        \frac{\norm{g_{k_i}}^2}{\mu_{k_i}}
         & \MPCchanged{\geq \frac{1}{\theta_{\max}} \frac{\norm{g_{k_i}}^2}{\sqrt{S_i}}} &  & (\text{definition}) \\
         & \MPCchanged{\geq \frac{1}{\theta_{\max}} \qty(\sqrt{S_i} - \sqrt{S_{i-1}})}. &  & (x/\sqrt{y} \geq \sqrt{y} - \sqrt{y-x})
    \end{align*}
    Summing these inequalities over $0 \le i < \abs{K^+}$ yields the claim.
    \myQED
\end{proof}

\section{Proof of \texorpdfstring{\cref{prop:bounded_BFGS}}{Proposition \ref{prop:bounded_BFGS}}}\label{app:proof_of_bounded_BFGS}

\MPCchanged{The proof of \cref{prop:bounded_BFGS} uses the standard BFGS update.}
\begin{proof}
    \MPCchanged{From \cref{eq:curvature_conditions1}, the initial scalar satisfies $0<\gamma=\norm{y_{k_{p-1}}}^2/(y_{k_{p-1}}^\top s_{k_{p-1}})\leq\Lambda$.}
    \MPCchanged{Let $B$ and $B_+$ denote the positive definite Hessian approximations before and after one update.}
    The BFGS update rule is
    \begin{equation*}
        \MPCchanged{B_+=B - \frac{Bss^\top B}{s^\top Bs} + \frac{yy^\top}{y^\top s}}.
    \end{equation*}
    \MPCchanged{Positive definiteness follows from $B\succ0$, $s\neq0$, and $y^\top s>0$.}
    \MPCchanged{Moreover, \cref{eq:curvature_conditions1} and the fact that $Bss^\top B/(s^\top Bs)$ is positive semidefinite give}
    \begin{equation*}
        \MPCchanged{B_+\preceq B+\frac{yy^\top}{y^\top s}\preceq B+\Lambda I}.
    \end{equation*}
    \MPCchanged{Applying this inequality recursively over $p$ updates from $B^{(0)}=\gamma I\preceq\Lambda I$ proves \cref{eq:bounded_BFGS_mM}.}
    \myQED
\end{proof}

\begin{MPCchange}
The fixed-initial-scale upper-curvature inequality is enforced directly by \cref{eq:implemented_cautious_test}, while the fixed shift in \cref{eq:proposed_choice_of_Bk} supplies the lower spectral bound.
No convexity or cocoercivity assumption is used.
This point is important because the convergence theorem permits nonconvex objectives.
\end{MPCchange}

\end{appendices}

\end{document}


\begin{appendices}
    \crefalias{section}{appendix}
    \section{Proof of \texorpdfstring{\cref{lem:sum_norm_g_over_mu_geq_2sqrt_leq_log}}{Lemma \ref{lem:sum_norm_g_over_mu_geq_2sqrt_leq_log}}}\label{sec:proof_of_sum_norm_g_over_mu_geq_2sqrt_leq_log}

The proof of \cref{lem:sum_norm_g_over_mu_geq_2sqrt_leq_log} is \MPCchanged{similar to previous arguments~\citep[Lemma~3.2]{wardAdaGradStepsizesSharp2020} and \citep[Lemma~3.1]{Gratton08022024}}.
\begin{proof}
    Recall that \cref{eq:offo_based_update} defines $\mu_k$ as $\theta_k \sqrt{\varsigma + \sum_{j \in K^+, \, j \leq k} \norm{g_j}^2}$, and \refLine{line:update_theta} of \cref{alg:regularized_qn} asserts $\theta_k \in [\theta_{\min}, \theta_{\max}]$.
    Let $k_i \in K^+$ be the $i$-th element of $K^+$.
    \MPCchanged{Define $S_{-1} \coloneqq \varsigma$ and $S_i \coloneqq \varsigma + \sum_{\ell=0}^{i} \norm{g_{k_\ell}}^2$ for $i \geq 0$.}
    For any $0 < x < y$, $x/y \le -\ln(1-x/y) = \ln y - \ln(y-x)$ holds.
    Thus, we have
    \begin{align*}
        \frac{\norm{g_{k_i}}^2}{\mu_{k_i}^2}
         & \MPCchanged{\leq \frac{1}{\theta_{\min}^2} \frac{\norm{g_{k_i}}^2}{S_i}} &  & (\text{definition}) \\
         & \MPCchanged{\leq \frac{1}{\theta_{\min}^2} \qty(\ln S_i - \ln S_{i-1})}.  &  & (x/y \leq \ln y - \ln(y-x))
    \end{align*}
    Similarly, $x/\sqrt{y} \ge x/(\sqrt{y}+\sqrt{y-x}) = \sqrt{y}-\sqrt{y-x}$ holds.
    Thus, we have
    \begin{align*}
        \frac{\norm{g_{k_i}}^2}{\mu_{k_i}}
         & \MPCchanged{\geq \frac{1}{\theta_{\max}} \frac{\norm{g_{k_i}}^2}{\sqrt{S_i}}} &  & (\text{definition}) \\
         & \MPCchanged{\geq \frac{1}{\theta_{\max}} \qty(\sqrt{S_i} - \sqrt{S_{i-1}})}. &  & (x/\sqrt{y} \geq \sqrt{y} - \sqrt{y-x})
    \end{align*}
    Summing these inequalities over $0 \le i < \abs{K^+}$ yields the claim.
    \myQED
\end{proof}

\section{Proof of \texorpdfstring{\cref{prop:bounded_BFGS}}{Proposition \ref{prop:bounded_BFGS}}}\label{app:proof_of_bounded_BFGS}

\MPCchanged{The proof of \cref{prop:bounded_BFGS} uses the standard BFGS update.}
\begin{proof}
    \MPCchanged{From \cref{eq:curvature_conditions1}, the initial scalar satisfies $0<\gamma=\norm{y_{k_{p-1}}}^2/(y_{k_{p-1}}^\top s_{k_{p-1}})\leq\Lambda$.}
    \MPCchanged{Let $B$ and $B_+$ denote the positive definite Hessian approximations before and after one update.}
    The BFGS update rule is
    \begin{equation*}
        \MPCchanged{B_+=B - \frac{Bss^\top B}{s^\top Bs} + \frac{yy^\top}{y^\top s}}.
    \end{equation*}
    \MPCchanged{Positive definiteness follows from $B\succ0$, $s\neq0$, and $y^\top s>0$.}
    \MPCchanged{Moreover, \cref{eq:curvature_conditions1} and the fact that $Bss^\top B/(s^\top Bs)$ is positive semidefinite give}
    \begin{equation*}
        \MPCchanged{B_+\preceq B+\frac{yy^\top}{y^\top s}\preceq B+\Lambda I}.
    \end{equation*}
    \MPCchanged{Applying this inequality recursively over $p$ updates from $B^{(0)}=\gamma I\preceq\Lambda I$ proves \cref{eq:bounded_BFGS_mM}.}
    \myQED
\end{proof}

\begin{MPCchange}
The fixed-initial-scale upper-curvature inequality is enforced directly by \cref{eq:implemented_cautious_test}, while the fixed shift in \cref{eq:proposed_choice_of_Bk} supplies the lower spectral bound.
No convexity or cocoercivity assumption is used.
This point is important because the convergence theorem permits nonconvex objectives.
\end{MPCchange}

\end{appendices}

\end{document}


\begin{appendices}
    \crefalias{section}{appendix}
    \section{Proof of \texorpdfstring{\cref{lem:sum_norm_g_over_mu_geq_2sqrt_leq_log}}{Lemma \ref{lem:sum_norm_g_over_mu_geq_2sqrt_leq_log}}}\label{sec:proof_of_sum_norm_g_over_mu_geq_2sqrt_leq_log}

The proof of \cref{lem:sum_norm_g_over_mu_geq_2sqrt_leq_log} is \MPCchanged{similar to previous arguments~\citep[Lemma~3.2]{wardAdaGradStepsizesSharp2020} and \citep[Lemma~3.1]{Gratton08022024}}.
\begin{proof}
    Recall that \cref{eq:offo_based_update} defines $\mu_k$ as $\theta_k \sqrt{\varsigma + \sum_{j \in K^+, \, j \leq k} \norm{g_j}^2}$, and \refLine{line:update_theta} of \cref{alg:regularized_qn} asserts $\theta_k \in [\theta_{\min}, \theta_{\max}]$.
    Let $k_i \in K^+$ be the $i$-th element of $K^+$.
    \MPCchanged{Define $S_{-1} \coloneqq \varsigma$ and $S_i \coloneqq \varsigma + \sum_{\ell=0}^{i} \norm{g_{k_\ell}}^2$ for $i \geq 0$.}
    For any $0 < x < y$, $x/y \le -\ln(1-x/y) = \ln y - \ln(y-x)$ holds.
    Thus, we have
    \begin{align*}
        \frac{\norm{g_{k_i}}^2}{\mu_{k_i}^2}
         & \MPCchanged{\leq \frac{1}{\theta_{\min}^2} \frac{\norm{g_{k_i}}^2}{S_i}} &  & (\text{definition}) \\
         & \MPCchanged{\leq \frac{1}{\theta_{\min}^2} \qty(\ln S_i - \ln S_{i-1})}.  &  & (x/y \leq \ln y - \ln(y-x))
    \end{align*}
    Similarly, $x/\sqrt{y} \ge x/(\sqrt{y}+\sqrt{y-x}) = \sqrt{y}-\sqrt{y-x}$ holds.
    Thus, we have
    \begin{align*}
        \frac{\norm{g_{k_i}}^2}{\mu_{k_i}}
         & \MPCchanged{\geq \frac{1}{\theta_{\max}} \frac{\norm{g_{k_i}}^2}{\sqrt{S_i}}} &  & (\text{definition}) \\
         & \MPCchanged{\geq \frac{1}{\theta_{\max}} \qty(\sqrt{S_i} - \sqrt{S_{i-1}})}. &  & (x/\sqrt{y} \geq \sqrt{y} - \sqrt{y-x})
    \end{align*}
    Summing these inequalities over $0 \le i < \abs{K^+}$ yields the claim.
    \myQED
\end{proof}

\section{Proof of \texorpdfstring{\cref{prop:bounded_BFGS}}{Proposition \ref{prop:bounded_BFGS}}}\label{app:proof_of_bounded_BFGS}

\MPCchanged{The proof of \cref{prop:bounded_BFGS} uses the standard BFGS update.}
\begin{proof}
    \MPCchanged{From \cref{eq:curvature_conditions1}, the initial scalar satisfies $0<\gamma=\norm{y_{k_{p-1}}}^2/(y_{k_{p-1}}^\top s_{k_{p-1}})\leq\Lambda$.}
    \MPCchanged{Let $B$ and $B_+$ denote the positive definite Hessian approximations before and after one update.}
    The BFGS update rule is
    \begin{equation*}
        \MPCchanged{B_+=B - \frac{Bss^\top B}{s^\top Bs} + \frac{yy^\top}{y^\top s}}.
    \end{equation*}
    \MPCchanged{Positive definiteness follows from $B\succ0$, $s\neq0$, and $y^\top s>0$.}
    \MPCchanged{Moreover, \cref{eq:curvature_conditions1} and the fact that $Bss^\top B/(s^\top Bs)$ is positive semidefinite give}
    \begin{equation*}
        \MPCchanged{B_+\preceq B+\frac{yy^\top}{y^\top s}\preceq B+\Lambda I}.
    \end{equation*}
    \MPCchanged{Applying this inequality recursively over $p$ updates from $B^{(0)}=\gamma I\preceq\Lambda I$ proves \cref{eq:bounded_BFGS_mM}.}
    \myQED
\end{proof}

\begin{MPCchange}
The fixed-initial-scale upper-curvature inequality is enforced directly by \cref{eq:implemented_cautious_test}, while the fixed shift in \cref{eq:proposed_choice_of_Bk} supplies the lower spectral bound.
No convexity or cocoercivity assumption is used.
This point is important because the convergence theorem permits nonconvex objectives.
\end{MPCchange}

\end{appendices}

\end{document}
